

\setcounter{section}{-1}


\section{Extended Materials and Methods}

\paragraph*{Model of Encoding and Decoding}
We assume a general encoding and decoding model for one-dimensional stimulus variables. Specifically, we assume that the stimulus space $\mathcal{X}$ is a circle or a bounded interval, the space of measurement (the neural encoding signal in the brain) is denoted as $\mathcal{Y}$, and the map $F : \mathcal{X} \rightarrow \mathcal{Y}$ is a bijection.
Formally, the neural encoding can be described by the following equation
\begin{equation}\label{eq:encoding}
    m = F(\theta) + \epsilon
\end{equation}
We assume, without loss of generality, that $\mathcal{Y}$ has volume 1.
Given a particular loss function $\ell$ and an encoding model $m\in\mathcal{Y}$, the Bayesian  estimator $\widehat{\theta} \in \mathcal{X}$ that minimizes the loss function is obtained as 
\begin{equation}\label{sec:lp-estimator}
	\widehat{\theta} := \arg_{\widehat\theta\in\mathcal{X}} \min \int_{\mathcal{X}}
	\ell({\widehat\theta}, \theta)
	p(\theta|m) d\theta
\end{equation}
At a stimulus $\theta$, the estimator $\widehat{\theta}$ has bias $\mathbb{E}[\widehat{\theta}]-\theta$.

The resource allocation is $\sqrt{\mathcal{J}(\theta)} = \frac{F'(\theta)}{\sigma}$, which is proportional to the slope $F'(\theta)$ of the encoding function,  with proportionality factor given by the reciprocal of the sensory noise SD. %
We study model variants where noise may also enter \emph{before} the map $F(\cdot)$ in SI Appendix S5.2.


%

\paragraph*{Formalizing the Space of Models}
As we consider the problem of identifying general observer models without fixed parametric forms for prior and encoding, the space of possible models is infinitely-dimensional. This idea can be conveniently formalized using standard  concepts from functional analysis.

We take $\mathcal{X}$ to be either a bounded interval or the unit circle. %
A function is continuous on the unit circle if it is continuous on $[0,2\pi)$ and additionally $f(2\pi) = f(0)$. 


Let $\mathcal{F}(\mathcal{X})$ be the set of real-valued functions on $\mathcal{X}$ that are five times continuously differentiable, whose values are always strictly positive, and where $\int f = 1$. Both the prior $\Prior$ and the resource allocation $F'(\theta)$ are elements of $\mathcal{F}(\mathcal{X})$.

We take $\mathcal{P} \subset \mathbb{N}$ to be a finite set of exponents $p$ for $L_p$ loss functions; we assume $\mathcal{P} \subset \{0,1,2,4,6,\dots\}$ for our theoretical results.

The space of encoding models is then the product space $\mathfrak{M} := \mathcal{F} \times \mathcal{F} \times \mathcal{P}$, whose elements are tuples of priors $\Prior$, resource allocations $F'$, and loss function exponents $p$.



We quantify volume on the infinite-dimensional space $\mathfrak{M}$ by %
parameterizing both $\Prior$ and $F'$ in terms of their logarithms, and placing a Gaussian process measure on these.
Independently sampling prior, encoding, and loss function determines a notion of measure or volume on the space of models.
See SI Appendix S2.1 for details.
%
%
%
The key here is that prior and encoding are sampled independently, so that there are no systematic algebraic relations between prior and encoding across the stimulus space. Our statement below that the exceptional set $\Omega$ has volume zero intuitively means that a sampled model is extremely unlikely to lie in $\Omega$.
%

%


\subsection*{Identifiability Theorems}
We consider the response distribution
$\mathbb{P}(y|M,\theta, \sigma)$
determined by the model $M \in \mathfrak{M}$, $\theta \in \mathcal{X}$, and a sensory noise magnitude $\sigma > 0$.
We assume that these distributions are known at each stimulus $\theta$, at one or more sensory noise levels.
The noise magnitudes $\sigma$ associated with the noise levels are assumed to be unknown, and inferred together with the other model components.
%


Our first theorem concerns recovery of the encoding, providing an explicit formula for the Fisher information in terms of the overlap in response distributions:



\begin{thm}[Recovering Encoding from Response Distribution]\label{thm:1}
    In the absence of motor noise, the encoding $F$ and the sensory noise magnitude can be recovered from the response distribution at a single level of sensory noise, with estimation error exponentially small in $\frac{1}{\sigma}$:
\begin{align*}
 \sqrt{\mathcal{J}(\theta)} = \sqrt{4\pi} \cdot \frac{\operatorname{d}}{\operatorname{d}h} \mathbb{P}\left(\widehat{\theta}(\theta+h) > \widehat{\theta}(\theta)\right) + O(\exp(-1/\sigma))
\end{align*}
where the derivative is evaluated at $h=0$, and $\widehat{\theta}(\theta+h)$ and $\widehat{\theta}(\theta)$ are independent response samples at stimuli $\theta+h$ and $\theta$, respectively.
%
\end{thm}
The proof of this theorem is in SI Appendix S2.4.
Numerical approximation of this formula allows straightforward estimation, independent of the loss function, as we illustrate in SI Appendix, Figure S3.



Our next theorem concerns recovery of the prior conditional on the loss function. It guarantees recovery of the prior when the loss function is known: %
\begin{thm}[One Level of Sensory Noise]\label{thm:2}
Let $\sigma > 0$ be the SD of sensory noise, and $p \in \mathcal{P}$.
Assume the %
response distribution $\mathbb{P}(\widehat{\theta}|M,\theta, \sigma)$ is given at each $\theta$, under a ground-truth model $\langle F', \Prior,p \rangle \in \mathfrak{M}$. 
There is a functional $\Phi_p$ mapping these distributions to $\langle \widehat{F}', \widehat{\Prior},p \rangle \in \mathfrak{M}$ such that 
\begin{align*}
|\widehat{F}(\theta) - F(\theta)| &= \mathcal{O}(\sigma^2) \\
|\widehat{\Prior}(\theta) - \Prior(\theta) | &= \mathcal{O}(\sigma^2) 
\end{align*}
as $\sigma \rightarrow 0$.
Constants in $\mathcal{O}(\cdot)$ depend on $F$ and $\Prior$, but not  $\theta$. 
\end{thm}
We next state our identifiability theorem for the more general setup where encoding, prior, and loss function are all unknown a priori.
We use $\Omega$ to denote the exceptional set of unidentifiable models.
\begin{thm}[Two Levels of Sensory Noise]\label{thm:3}
There is a subset $\Omega \subset \mathfrak{M}$ of volume $0$ such that the following holds.
Let $0 < \sigma_1 < \sigma_2$ be two levels of sensory noise. %
Assume the %
response distributions $\mathbb{P}(\widehat{\theta}|M,\theta, \sigma_i)$ are given at each $\theta$ and for both $\sigma_1, \sigma_2$, under a ground-truth model $M = \langle F', \Prior,p \rangle \in \mathfrak{M}$.
There is a functional $\Phi$ mapping the collection of these distributions to $\langle \widehat{F}', \widehat{\Prior},q \rangle \in \mathfrak{M}$
such that---provided $M \not\in \Omega$---we have
\begin{align*}
    |\widehat{F}(\theta) - F(\theta)| &= \mathcal{O}(\sigma_2^2) \\
    |\widehat{\Prior}(\theta) - \Prior(\theta)| &= \mathcal{O}(\sigma_2^2) \\
    q &= p  \text{ for $\sigma_2^2$ small}
\end{align*}
%
Constants in the $\mathcal{O}(\cdot)$ expressions depend on the ratio $\frac{\sigma_2}{\sigma_1}$, and on $F$ and $\Prior$, but not on $\theta$.

%
\end{thm}
%
This theorem provides in-principle identification of the loss function in all models outside of $\Omega$. %
Importantly, this theorem does not presuppose that the magnitudes $\sigma_1, \sigma_2$ are known. It only requires that $\sigma_1, \sigma_2$ are small and distinct.
Although in theory, two levels of noise is sufficient, practically, it may be advantageous to obtain observations at a larger number of noise levels across different scales according to our numerical results. 
%
%
%



\begin{thm}[Role of Motor Noise]\label{thm:4}
%


Theorem~\ref{thm:3} is not affected by adding symmetric isotropic motor noise of variance $\tau^2$, and by guessed responses appearing at a rate $0 < \gamma < 1$. The constants in the $\mathcal{O}(\cdot)$ expressions then depend additionally on the ratio $\frac{\tau}{\sigma_1}$ and on $\gamma$.
\end{thm}







The proofs of Theorems~\ref{thm:2}--\ref{thm:4}, provided in SI Appendix S2.5 and S2.6, recover model components using the Method of Moments, %
by matching observed moments of the response distribution to truncated Taylor expansions of the moments in powers of the noise magnitude $\sigma^2$, with coefficients reflecting the model components $F', \Prior, p$.
The key intuition behind the proofs of Theorems~\ref{thm:3}--\ref{thm:4} is that observing multiple sensory noise levels allows recovering multiple coefficients of these Taylor expansions. Even if the lowest-order coefficient does not allow uniquely identifying the model, the next-order coefficient provides sufficient information to uniquely identify any model outside of the exceptional set $\Omega$ of unidentifiable models.



This strategy largely recovers the ground-truth model up to an error indicating deviation from the low-noise regime.
Our simulation experiments indicate that conclusions remain valid at finite noise, and when only a finite sample dataset is available.




\subsection*{Implementation}
We use the fitting procedure introduced in \cite{hahn2024unifying}, which, given a loss function, jointly fits encoding, prior, motor and sensory noise magnitudes, and guessing rate, by maximizing the data's likelihood under the model. Prior and encoding are specified on a discretized grid over $\mathcal{X}$. Following \cite{hahn2024unifying}, its size is 180 in the models in Figures 1--7, and 200 in the model in Figure~\ref{fig:remington}. Following \cite{hahn2024unifying}, we implement the $L_p$ loss function on circular spaces via powers of the cosine ($\ell(x,y) := (1-\cos(x-y))^{p/2}$), as they found this to provide better fit to behavioral data than a version based on absolute distance. In the low-noise setting, $\ell(x,y) \sim |x-y|^p$, i.e., our theoretical results apply to this implementation.

We provide a description of the practical model fitting procedure in SI Appendix, Section S8. As in \cite{hahn2024unifying}, the fitting procedure includes regularization encouraging smooth fits for prior and encoding.
For synthetic data in Figures 2--7, we used the same regularization strength as \cite{hahn2024unifying} used for the data collected by \cite{deGardelle2010AnOI}.
In Figure~\ref{fig:remington}, we used the same regularization strength as in the fit to the original data from \cite{hahn2024unifying}.


Note that this method is applicable even outside of the small-noise regime, allowing us to verify the generalizability of the theoretical conclusions beyond it.



The accompanying code base provides functionality to (i) sample datasets from specified models, and (ii) fit models at different loss function exponents to these datasets using the fitting procedure described above.
The implementation uses Python 3.9, matplotlib 3.9.4, numpy 2.0.2, scipy 1.13.1, torch 2.5.1.
A demo using a synthetic dataset (1K trials) is available with the code base.



\subsection*{Simulations}
Following \cite{hahn2024unifying}, simulated datasets include Gaussian (or von Mises) motor noise and guesses. For all experiments on circular stimulus spaces, motor noise magnitude and guessing rate are taken from the fit to the data of \cite{deGardelle2010AnOI} as obtained by \cite{hahn2024unifying}. We also use sensory magnitudes from the same fit; focusing on the four noise levels corresponding to exposure durations of 40ms, 80ms, 160ms, 1000ms, as we find four levels to largely be sufficient for identifiability.
%

%

%
%
The numbers of trials in simulated data are as indicated in Fig.~\ref{fig:2}, 10K in Fig.~\ref{fig:confounded-prior}, 10K in Fig.~\ref{fig:random-models}, and as indicated in Fig.~\ref{fig:levels-sizes-synthetic}.

%
For experiments on circular stimulus spaces, we focus on the orientation domain, with stimulus space $[0^\circ, 180^\circ]$. 


%
As in \cite{hahn2024unifying}, we randomly partition each dataset into 10 folds. We fit models on 9 and evaluate negative log-likelihood (NLL) on one fold. We plot $\Delta$NLL relative to the NLL achieved by the ground-truth loss function exponent.

Confidence bands in Figure~\ref{fig:2} are obtained by resampling 2K replicates of the dataset, and obtaining 95\% confidence intervals for the values at 45 regularly spaced grid points with Bonferroni correction.

%
For Fig.~\ref{fig:random-models}, we randomized models as follows.
We sampled $F'$ and $\Prior$ each by parameterizing their logarithms in terms of the first four terms of their Fourier series, with coefficients uniformly in $[-0.5, 0.5]$.


%
For Fig.~\ref{fig:levels-sizes-synthetic}, we considered (i) uniform prior and periodic resources $\sqrt{\FI(\theta)} \propto \frac{3}{2} - \left|\sin(\frac{\theta}{2})\right|$ (stylized from the fit to the data collected by \cite{deGardelle2010AnOI}), (ii) matched periodic prior and encoding, (iii) periodic prior and uniform encoding, (iv) a periodic resource allocation and a shifted prior $\Prior \propto \frac{3}{2} - \left|\sin(\frac{\theta+\pi}{2})\right|$.
The same periodic encoding was used in Figs.~\ref{fig:2} and \ref{fig:confounded-prior}.




%
In Fig.~\ref{fig:recovered-fi}, we show results fitted at $p=8$, selected based on quality of fit on the full dataset.  Fit of the encoding is highly consistent even when fitting at other exponents (SI Appendix, Figure S33).
In Fig.~\ref{fig:recovered-fi}c, we obtained the environmental prior from the ``Blended'' distribution in Supplementary Figure 4 of \cite{Girshick2011CardinalRV}, at 640 x 480 resolution. We further obtained discrimination thresholds from Figure 1 in \cite{westheimer2003meridional}, averaging over the five subjects reported there.

%

%
In Fig.~\ref{fig:gardelle-fits}, we averaged the NLL over five random seeds for downsampling the dataset.



%
In Fig.~\ref{fig:remington}, parameters for the unimodal model are taken from the fit to the behavioral data, which is from \cite{hahn2024unifying}. Parameters for the bimodal model are identical except for the prior.
%
In Fig.~\ref{fig:remington}, all fits presuppose a logarithmic encoding consistent with Weber's law, as in the analysis of the dataset in \cite{hahn2024unifying}. Corresponding model fit statistics for freely fitted encodings (SI Appendix, Fig. S34) show that the bimodal model is identifiable even without assuming a specific encoding.

\subsection*{Experimental data} The experimental data used in this paper are published in ref \cite{deGardelle2010AnOI} and \cite{Remington2018LateBI}. Data collection and analysis were not performed blind to the conditions of the experiments.

%
Orientation perception \citep{deGardelle2010AnOI}: 49 subjects adjusted a blue strip in order to reproduce the orientation of a Gabor patch.
Stimuli were uniformly distributed on most trials, but had orientation 0, 45, 90, 135 degrees on a specific subset of trials; we followed \cite{hahn2024unifying} in excluding the latter.
Depending on the trial, exposure duration was 20, 40, 80, 160, or 1000 ms.
On 10\% of trials, no stimulus was shown; we excluded these.
There were 2,208 trials at 20ms,  40ms, 80ms, 160ms each, and 1,104 trials at 1000ms. In total, 9,936 trials were included.



Time intervals \citep{Remington2018LateBI}: 15 subjects reproduced time intervals in a Ready-Set-Go Task in Experiments 1 and 2.
Sessions differed in whether subjects reproduced the time interval itself (the Identity condition)  or reproduced it scaled by factors of 0.75 or 1.5; we only included data in the identity condition.  In total, 8,999 trials were included.


